\section{Empirical Results}
\label{sec:results}
% ============================================================

\subsection{Experimental Setup}
\label{sec:setup}

\textbf{Benchmark.}  We regenerated the architecture sweep for this version.  It contains 474 CIFAR-10 networks:
78 MLPs, 72 CNNs, 60 depthwise-separable CNNs, 132 ResNet-style and 132 depthwise ResNet-style networks,
each with ReLU, $x^2$, $x^3{+}x$ or $x^2{+}x$ activations (111, 121, 121 and 121 networks) and no normalization, LayerNorm or GroupNorm.
CIFAR-100 and ImageNet-16-120 results from an earlier version are withdrawn: the underlying files could not be recovered
and the CIFAR-100 ReLU subset had only three networks.

\textbf{Training.}  30-epoch SGD (momentum $0.9$, weight decay $5{\times}10^{-4}$, cosine schedule, batch $128$, random crop with padding $4$
and horizontal flip), learning rate $0.1$, or $0.01$ with gradient clipping $1.0$ for polynomial networks without normalization.  The sweep is
run by \texttt{sweep\_gpu.py}, which keeps this recipe but holds the dataset and the augmentation on the GPU.  As a check against the original CPU pipeline,
three architectures trained for the full $30$ epochs give $77.45\%$ for the CNN ReLU (earlier $77.1\%$; three-seed mean $76.98\%$), $75.42\%$ for the DWS ReLU
(earlier $75.0\%$; $74.73\%$) and $54.14\%$ for the $x^2{+}x$ MLP (earlier $51.1\%$; $50.67\%$); the last is about three points higher than before, so we do
not claim exact equivalence to the earlier pipeline.  Every architecture is trained with two seeds (0 and 1).  Zero-cost proxies and the cost estimate are computed
once, at the seed-0 initialization.  A network is called \emph{failed} if training diverged or its mean test accuracy is below $15\%$.

\textbf{ezkl.}  All calibrations in this section use ezkl 9.5.0 (\texttt{calibrate\_settings}, \texttt{target=resources}), $20$ real CIFAR-10 test images as calibration
inputs, and the seed-0 \emph{trained} weights, on a 24-logical-core AMD Threadripper PRO 3945WX host shared with other jobs.  Networks that failed to train
(nothing meaningful to calibrate) and networks whose estimated $\logrows$ exceeds $25$ (irrelevant to the budgets studied) were not calibrated; the list is released.
Of the 255 calibrations attempted, 12 failed inside ezkl and are excluded from pools.

\textbf{Released data.}  Per-architecture results for both seeds (accuracy, proxies at initialization, cost estimate), the calibration records with the per-operator
features used for fitting, and every summary number in this section are in the repository, together with the scripts that regenerate Tables~\ref{tab:corr}--\ref{tab:ezkl}.
The raw output of the earlier three-seed job (Section~\ref{sec:robust}) is not released.

\subsection{Label Noise and Training Failures}

Across the two seeds the median per-architecture standard deviation of test accuracy is 0.24 points (90th percentile 1.53; the maximum, 45.4,
is a network that trained in one seed and failed in the other), and the Spearman correlation between the two seeds' accuracies is 0.95.
Label noise is therefore small relative to differences between architectures.  \emph{Training failure} is large and depends on the recipe:
153 of the 474 networks have mean accuracy below $15\%$.  Failures by activation are 0 of 111 (ReLU), 50 of 121 ($x^2$),
73 of 121 ($x^3{+}x$) and 47 of 121 ($x^2{+}x$); by normalization, 37 none, 49 LayerNorm and 84 GroupNorm
(polynomial networks with normalization use the larger learning rate, so normalization does not prevent failure here).  The accuracy labels of polynomial networks are therefore
partly trainability labels, and we report results with and without the failed networks.

\subsection{Earlier Multi-Seed and Training-Budget Validation}
\label{sec:robust}

Table~\ref{tab:robustness} reports the replication of the eight architectures that
appeared in the earlier version's timing table.  All values are test top-1 accuracy; multi-seed entries are
mean $\mathbin{\pm}$ sample standard deviation over three seeds.  Every
single-seed value is within $1.2$ percentage points of the corresponding
three-seed 30-epoch mean.  In particular, the DWS ReLU (32/64/128/256) network, the architecture selected in the earlier version's search, replicates at $74.73\%\mathbin{\pm}1.61\%$ and reaches
$79.96\%\mathbin{\pm}0.40\%$ after 100 epochs.

The activation-matched controls reinforce that activation effects are
architecture-dependent.  ReLU remains the highest-accuracy choice in each
matched convolutional group, while $x^2{+}x$ is the strongest polynomial
alternative.  In the normalization-free MLP, however, $x^2{+}x$ reaches
$55.67\%$ at 100 epochs versus $53.52\%$ for ReLU.  Pure-$x^2$ DWS controls
remain at chance accuracy under both budgets, so polynomial activations should
not be treated as interchangeable.

\begin{table}[h]
\centering
\caption{Multi-seed replication and longer-training validation of the eight
  architectures used in the accuracy/proving table.  ``Original'' is the
  reported single-seed result; validation columns use fresh training runs.}
\label{tab:robustness}
\small
\begin{tabular}{lrrr}
\toprule
\textbf{Architecture} & \textbf{Original 30 ep.} & \textbf{3-seed 30 ep.} & \textbf{3-seed 100 ep.} \\
\midrule
MLP $x^2{+}x$, none (128/64/32)       & $51.1$ & $50.67\mathbin{\pm}0.52$ & $55.67\mathbin{\pm}0.32$ \\
MLP $x^2{+}x$, LN (128/64/32)         & $52.0$ & $50.87\mathbin{\pm}1.16$ & $54.10\mathbin{\pm}1.77$ \\
DWS ReLU (16/32/64)                   & $62.9$ & $64.10\mathbin{\pm}0.52$ & $67.52\mathbin{\pm}1.12$ \\
DWS ReLU (32/64/128)                  & $69.3$ & $70.48\mathbin{\pm}1.30$ & $75.70\mathbin{\pm}0.71$ \\
DWS ReLU (32/64/128/256)              & $75.0$ & $74.73\mathbin{\pm}1.61$ & $79.96\mathbin{\pm}0.40$ \\
CNN ReLU (16/32/64)                   & $77.1$ & $76.98\mathbin{\pm}0.72$ & $79.80\mathbin{\pm}0.46$ \\
CNN ReLU (16/32/64/128)               & $81.6$ & $81.31\mathbin{\pm}0.53$ & $84.20\mathbin{\pm}0.32$ \\
CNN ReLU (64/128/256)                 & $83.9$ & $83.96\mathbin{\pm}0.16$ & $86.25\mathbin{\pm}0.21$ \\
\bottomrule
\end{tabular}
\end{table}

Across all 26 architecture means, the 30- to 100-epoch Spearman correlation
is $\rho=0.932$.  Across the 77 successful architecture-seed pairs it is
$\rho=0.941$, and the three seed-wise correlations range from $0.929$ to
$0.964$.  Absolute accuracy rises by $5.88$ percentage points on average, so
the longer budget changes accuracy levels more than relative ordering.  The
ordering is strong but not perfect: some deep polynomial CNNs learn late.
Of 156 planned runs, 154 complete; both divergences are the same
$x^3{+}x$ LayerNorm MLP under seed 44, once at each epoch budget, with a
non-finite loss during epoch 1.


\subsection{Zero-Cost Proxies Across Activation Families}

Table~\ref{tab:corr} gives Spearman $\rho$ between each proxy and mean accuracy per activation family with bootstrap intervals and Holm-adjusted
significance over the 32 activation-family tests it contains.  NEAR is a reliable predictor only for ReLU networks
($\rho=+0.70$, $95\%$ CI $[+0.57,+0.80]$); on the polynomial families its intervals include zero
($x^2$ -0.12 $[-0.29,+0.05]$, $x^3{+}x$ +0.01, $x^2{+}x$ -0.14).  ZenNAS is negatively correlated in every family
(ReLU -0.23, $x^2$ -0.42, $x^3{+}x$ -0.31, $x^2{+}x$ -0.54).  ZiCo is the one proxy with a significant positive coefficient on a polynomial family
($x^3{+}x$, +0.45), but once failed networks are removed (Table~\ref{tab:corrnofail}) it is -0.38 on $x^3{+}x$ and -0.50 on $x^2$:
its apparent success on $x^3{+}x$ comes from separating networks that train from networks that do not, not from ranking accuracy among those that train.
MeCo and Synflow are weak or inconsistent.

Two controls address confounds.  Table~\ref{tab:corrsame} restricts to normalized networks, where every activation was trained with the same recipe, and the picture is the same
(NEAR on ReLU +0.73, on $x^2$ -0.02, $x^3{+}x$ +0.06, $x^2{+}x$ -0.07).  Table~\ref{tab:partial} removes, by rank regression, the effects of architecture type, normalization and
parameter count: NEAR's partial correlation on ReLU is +0.22 (against +0.70 unadjusted), so much of NEAR's apparent power on ReLU networks reflects architecture type and
size.  We therefore do not claim that any tested proxy is a reliable accuracy predictor across activation families, and we do not attribute the failure to the activation function alone.

\begin{table}[h]
\centering
\caption{Spearman $\rho$ between proxy score at initialization and mean top-1 accuracy over two training seeds
  (CIFAR-10, regenerated sweep), per activation family, with $95\%$ bootstrap CI over architectures.
  ${}^{*}$ marks Holm-adjusted $p<0.05$ over the 25 activation-family tests in this table.
  A dash or ``n/a'' means fewer than $8$ finite scores.}
\label{tab:corr}
\small
\begin{tabular}{lrrrr}
\toprule
\textbf{Proxy} & \textbf{ReLU} ($n{=}111$) & \textbf{$x^2$} ($n{=}121$) & \textbf{$x^3{+}x$} ($n{=}121$) & \textbf{$x^2{+}x$} ($n{=}121$) \\
\midrule
NEAR & $+0.70^{*}$ {\scriptsize[+0.57,+0.80]} & $-0.12$ {\scriptsize[-0.29,+0.05]} & $+0.01$ {\scriptsize[-0.17,+0.19]} & $-0.14$ {\scriptsize[-0.30,+0.03]} \\
Synflow & $+0.22$ {\scriptsize[+0.03,+0.40]} & $-0.24$ {\scriptsize[-0.41,-0.04]} {\scriptsize($n{=}99$)} & $+0.13$ {\scriptsize[-0.14,+0.43]} {\scriptsize($n{=}70$)} & $-0.18$ {\scriptsize[-0.38,+0.03]} {\scriptsize($n{=}99$)} \\
ZiCo & $+0.24$ {\scriptsize[+0.07,+0.40]} & $-0.10$ {\scriptsize[-0.29,+0.09]} & $+0.45^{*}$ {\scriptsize[+0.29,+0.60]} & $-0.04$ {\scriptsize[-0.22,+0.16]} \\
MeCo & $+0.16$ {\scriptsize[-0.02,+0.32]} & $-0.09$ {\scriptsize[-0.26,+0.08]} {\scriptsize($n{=}109$)} & $+0.05$ {\scriptsize[-0.16,+0.27]} {\scriptsize($n{=}79$)} & $-0.06$ {\scriptsize[-0.24,+0.14]} {\scriptsize($n{=}109$)} \\
Zen-style (variant) & $-0.23$ {\scriptsize[-0.42,-0.04]} & $-0.42^{*}$ {\scriptsize[-0.54,-0.27]} {\scriptsize($n{=}109$)} & $-0.31$ {\scriptsize[-0.49,-0.11]} {\scriptsize($n{=}79$)} & $-0.54^{*}$ {\scriptsize[-0.64,-0.42]} {\scriptsize($n{=}109$)} \\
NASWOT (ReLU patterns) & $+0.60^{*}$ {\scriptsize[+0.44,+0.75]} & n/a & n/a & n/a \\
NASWOT (pre-activation signs) & $+0.60^{*}$ {\scriptsize[+0.44,+0.75]} & $-0.27^{*}$ {\scriptsize[-0.43,-0.11]} & $-0.29^{*}$ {\scriptsize[-0.45,-0.12]} & $-0.35^{*}$ {\scriptsize[-0.49,-0.19]} \\
\bottomrule
\end{tabular}
\end{table}

\begin{table}[h]
\centering
\caption{As Table~\ref{tab:corr} after excluding networks whose training diverged or ended at chance level
  (mean top-1 $<15\%$; 0 of 111 ReLU, 50 of 121 $x^2$, 73 of 121 $x^3{+}x$, 47 of 121 $x^2{+}x$ excluded).  Holm adjustment is applied within this table.}
\label{tab:corrnofail}
\small
\begin{tabular}{lrrrr}
\toprule
\textbf{Proxy} & \textbf{ReLU} ($n{=}111$) & \textbf{$x^2$} ($n{=}71$) & \textbf{$x^3{+}x$} ($n{=}48$) & \textbf{$x^2{+}x$} ($n{=}74$) \\
\midrule
NEAR & $+0.70^{*}$ {\scriptsize[+0.57,+0.80]} & $-0.16$ {\scriptsize[-0.38,+0.09]} & $+0.09$ {\scriptsize[-0.21,+0.39]} & $-0.07$ {\scriptsize[-0.31,+0.20]} \\
Synflow & $+0.22$ {\scriptsize[+0.03,+0.40]} & $-0.49^{*}$ {\scriptsize[-0.64,-0.27]} {\scriptsize($n{=}66$)} & $-0.35$ {\scriptsize[-0.62,+0.01]} {\scriptsize($n{=}40$)} & $-0.38^{*}$ {\scriptsize[-0.58,-0.15]} {\scriptsize($n{=}69$)} \\
ZiCo & $+0.24$ {\scriptsize[+0.07,+0.40]} & $-0.50^{*}$ {\scriptsize[-0.69,-0.26]} & $-0.38$ {\scriptsize[-0.60,-0.12]} & $-0.37^{*}$ {\scriptsize[-0.58,-0.12]} \\
MeCo & $+0.16$ {\scriptsize[-0.02,+0.32]} & $-0.20$ {\scriptsize[-0.42,+0.04]} & $-0.03$ {\scriptsize[-0.33,+0.26]} & $-0.15$ {\scriptsize[-0.40,+0.11]} \\
Zen-style (variant) & $-0.23$ {\scriptsize[-0.42,-0.04]} & $-0.18$ {\scriptsize[-0.42,+0.12]} & $-0.49^{*}$ {\scriptsize[-0.70,-0.20]} & $-0.39^{*}$ {\scriptsize[-0.59,-0.12]} \\
NASWOT (ReLU patterns) & $+0.60^{*}$ {\scriptsize[+0.44,+0.75]} & n/a & n/a & n/a \\
NASWOT (pre-activation signs) & $+0.60^{*}$ {\scriptsize[+0.44,+0.75]} & $-0.25$ {\scriptsize[-0.46,+0.00]} & $-0.26$ {\scriptsize[-0.51,+0.05]} & $-0.26$ {\scriptsize[-0.47,-0.02]} \\
\bottomrule
\end{tabular}
\end{table}

\begin{table}[h]
\centering
\caption{As Table~\ref{tab:corr}, restricted to networks with LayerNorm or GroupNorm.  All activations in this subset were trained with the same
  recipe (lr $0.1$), removing the learning-rate difference that applies to norm-free polynomial networks (lr $0.01$ with clipping).
  Holm adjustment is applied within this table.}
\label{tab:corrsame}
\small
\begin{tabular}{lrrrr}
\toprule
\textbf{Proxy} & \textbf{ReLU} ($n{=}72$) & \textbf{$x^2$} ($n{=}77$) & \textbf{$x^3{+}x$} ($n{=}77$) & \textbf{$x^2{+}x$} ($n{=}77$) \\
\midrule
NEAR & $+0.73^{*}$ {\scriptsize[+0.58,+0.82]} & $-0.02$ {\scriptsize[-0.23,+0.18]} & $+0.06$ {\scriptsize[-0.17,+0.29]} & $-0.07$ {\scriptsize[-0.27,+0.13]} \\
Synflow & $+0.19$ {\scriptsize[-0.09,+0.47]} & $-0.52^{*}$ {\scriptsize[-0.65,-0.34]} {\scriptsize($n{=}69$)} & $-0.29$ {\scriptsize[-0.54,+0.01]} {\scriptsize($n{=}49$)} & $-0.58^{*}$ {\scriptsize[-0.70,-0.40]} {\scriptsize($n{=}69$)} \\
ZiCo & $+0.32$ {\scriptsize[+0.08,+0.51]} & $-0.09$ {\scriptsize[-0.30,+0.14]} & $+0.42^{*}$ {\scriptsize[+0.22,+0.60]} & $-0.03$ {\scriptsize[-0.25,+0.18]} \\
MeCo & $+0.17$ {\scriptsize[-0.05,+0.38]} & $-0.13$ {\scriptsize[-0.34,+0.10]} {\scriptsize($n{=}69$)} & $+0.17$ {\scriptsize[-0.10,+0.41]} {\scriptsize($n{=}49$)} & $+0.04$ {\scriptsize[-0.19,+0.28]} {\scriptsize($n{=}69$)} \\
Zen-style (variant) & $+0.19$ {\scriptsize[-0.07,+0.44]} & $-0.53^{*}$ {\scriptsize[-0.67,-0.34]} {\scriptsize($n{=}69$)} & $-0.26$ {\scriptsize[-0.52,+0.07]} {\scriptsize($n{=}49$)} & $-0.51^{*}$ {\scriptsize[-0.65,-0.32]} {\scriptsize($n{=}69$)} \\
NASWOT (ReLU patterns) & $+0.60^{*}$ {\scriptsize[+0.38,+0.75]} & n/a & n/a & n/a \\
NASWOT (pre-activation signs) & $+0.59^{*}$ {\scriptsize[+0.38,+0.75]} & $-0.24$ {\scriptsize[-0.44,-0.03]} & $-0.08$ {\scriptsize[-0.31,+0.14]} & $-0.31$ {\scriptsize[-0.48,-0.12]} \\
\bottomrule
\end{tabular}
\end{table}

\begin{table}[h]
\centering
\caption{Partial Spearman correlation of proxy and accuracy after removing, by rank regression, the effects of
  architecture type (MLP, CNN, DWS-CNN, ResNet-style), normalization and $\log$ parameter count.  ${}^{*}$: unadjusted
  $p<0.05$ for the partial correlation (descriptive only; the Holm-adjusted tests are in Table~\ref{tab:corr}).}
\label{tab:partial}
\small
\begin{tabular}{lrrrr}
\toprule
\textbf{Proxy} & \textbf{ReLU} & $\bm{x^2}$ & $\bm{x^3{+}x}$ & $\bm{x^2{+}x}$ \\
\midrule
NEAR & $+0.22^{*}$ & $+0.05$ & $+0.03$ & $-0.01$ \\
Synflow & $+0.40^{*}$ & $-0.37^{*}$ & $-0.19$ & $-0.31^{*}$ \\
ZiCo & $+0.13$ & $-0.06$ & $+0.48^{*}$ & $+0.01$ \\
MeCo & $+0.23^{*}$ & $-0.05$ & $-0.01$ & $-0.01$ \\
Zen-style (variant) & $+0.02$ & $-0.32^{*}$ & $-0.12$ & $-0.36^{*}$ \\
NASWOT (ReLU patterns) & $+0.00$ & -- & -- & -- \\
NASWOT (pre-activation signs) & $-0.00$ & $-0.07$ & $-0.06$ & $-0.20^{*}$ \\
\bottomrule
\end{tabular}
\end{table}


\subsection{NASWOT}

Standard NASWOT scores the outputs of \texttt{nn.ReLU} modules.  It returns $-\infty$ on all 363 of the 363 polynomial networks because they contain none; for the $x^2$ networks the degeneracy
is also mathematically correct (Proposition~\ref{prop:naswot}), but the count does not test the proposition.  To remove the confound we score the sign of each activation module's input instead
(Section~\ref{sec:theory}).  That variant equals the standard score on ReLU networks ($\rho=+0.60$) and is defined for every polynomial network, yet it is negatively correlated with accuracy there
($x^2$ -0.27, $x^3{+}x$ -0.29, $x^2{+}x$ -0.35; Table~\ref{tab:corr}).  The failure of NASWOT on polynomial networks is
therefore not only a missing-module artifact.

\begin{table}[h]
\centering
\caption{Networks with a $-\infty$ NASWOT score: standard implementation (scores \texttt{nn.ReLU} outputs) versus the
  pre-activation-sign variant.}
\label{tab:naswot}
\small
\begin{tabular}{lrrr}
\toprule
\textbf{Family} & $n$ & \textbf{standard $-\infty$} & \textbf{sign variant $-\infty$} \\
\midrule
ReLU & 111 & 0 & 0 \\
$x^2$ & 121 & 121 & 0 \\
$x^3{+}x$ & 121 & 121 & 0 \\
$x^2{+}x$ & 121 & 121 & 0 \\
\bottomrule
\end{tabular}
\end{table}


\subsection{ZenNAS}

Among the 110 polynomial networks without normalization, the Zen score is weakly negatively correlated with accuracy ($\rho=-0.22$), positively correlated with depth ($+0.35$)
and with $\log$ parameter count ($+0.15$), and the failure rate is 43\% in the top Zen quartile against 11\% in the bottom quartile
(14\% overall).  Excluding failed networks, the correlation with accuracy is -0.19 ($n=95$).  Among the 39 ReLU networks without normalization it is -0.42, with $\log$ parameter count +0.40.
This is descriptive: it is consistent with Zen rewarding size and depth, which are not what determines accuracy here, but we do not claim a mechanism.

\subsection{ezkl Circuit Size and Timings on Trained Networks}
\label{sec:ezkltrained}

\textbf{Width-matched comparison.}  Table~\ref{tab:matchedtrained} repeats the width-matched comparison of Section~\ref{sec:findings} on trained networks calibrated with real inputs.
The result is unchanged: ReLU and $x^2$ networks have the same $\logrows$ in all five unnormalized skeletons, and unnormalized $x^3{+}x$ networks calibrate to $\logrows{=}23$ in every skeleton.
Across all 52 ReLU/$x^2$ pairs in the sweep that share a skeleton and were calibrated, 47 have the same $\logrows$, in 4 the $x^2$ network is smaller and in 1 it is larger; the median
ratio of assigned cells ($x^2$/ReLU) is 1.03.

\textbf{Accounting.}  $\lceil\log_2 C\rceil$ equals the calibrated $\logrows$ for 153/201 trained ReLU, $x^2$ and $x^2{+}x$ networks (ReLU 62/80, $x^2$ 44/57, $x^2{+}x$ 47/64) and for
19/42 of the $x^3{+}x$ networks.  The identity is therefore not exact for trained networks with normalization or residual structure, where calibration sometimes enlarges the circuit.

\textbf{Timings.}  Table~\ref{tab:ezkl} replaces the timings of the earlier version, whose ezkl version, settings, weight state and hardware were not recorded.  It covers 7 trained seed-0 networks at $\logrows$ 19--22
with real inputs, one ezkl version and one host.  Mock proofs were repeated three times (spread at most 19\% of the median) and take 2--42\,s.
Real proofs were run for the 4 networks up to $\logrows{=}21$, three times when $\logrows<20$ and once otherwise, and take 25--350\,s; larger circuits were mock-proved only.
The proving times in the earlier version (up to $1340$\,s for a $\logrows{=}21$ network) are \emph{not} reproduced: the comparable network here proves in a fraction of that time, and we withdraw the earlier timings.
The quantized circuits agree with the float models on 700 of 700 real test-image predictions.  The host is shared with other jobs, so times are indicative; the load average before and after each run is in the released file.

\begin{table}[h]
\centering
\caption{Width-matched ezkl circuit size on \emph{trained} seed-0 networks calibrated with real CIFAR-10 inputs: assigned cells (calibrated $\logrows$).
  ``fail'': calibration did not complete or errored (recorded in the released file).}
\label{tab:matchedtrained}
\small
\begin{tabular}{lrrrr}
\toprule
\textbf{Skeleton} & \textbf{ReLU} & $\mathbf{x^2}$ & $\mathbf{x^2{+}x}$ & $\mathbf{x^3{+}x}$ \\
\midrule
MLP 128-64-32 & 0.404M (19) & 0.407M (19) & 0.408M (19) & 0.402M (23) \\
MLP 256-128 & 0.823M (20) & 0.824M (20) & 0.824M (20) & 0.813M (23) \\
MLP 512-256 & 1.703M (21) & 1.710M (21) & 1.711M (21) & 1.688M (23) \\
CNN 16-32 & 1.732M (21) & 1.832M (21) & 1.855M (21) & 1.879M (23) \\
CNN 16-32-64 & 2.795M (22) & 2.909M (22) & 2.938M (22) & 2.956M (23) \\
\bottomrule
\end{tabular}
\end{table}

\begin{table}[h]
\centering
\caption{ezkl timings on trained seed-0 networks with real CIFAR-10 inputs under one recorded configuration
  (ezkl 9.5.0, AMD Ryzen Threadripper PRO 3945WX 12-Cores, 24 logical cores, shared host).  Median [min,max] over
  repeats; seconds.  Mock proofs: $\times3$.  Real proofs: $\times3$ for $\logrows\le21$, $\times1$ for $\logrows=22$.  Last column: quantized-circuit predictions that agree with the float model on real test images.}
\label{tab:ezkl}
\scriptsize
\begin{tabular}{lrrrrrrr}
\toprule
\textbf{Architecture} & \textbf{top-1 (2 seeds)} & $\logrows$ & \textbf{mock} & \textbf{prove} & \textbf{verify} & \textbf{proof KB} & \textbf{q-agree} \\
\midrule
mlp\_x2px\_none\_128x64x32 & 54.5 & 19 & 2.5 {\scriptsize[2.5,2.5]} & 28.3 {\scriptsize[28.2,28.6]} & 0.20 {\scriptsize[0.20,0.21]} & 15.4 & 100/100 \\
mlp\_x2px\_layernorm\_128x64x32 & 56.5 & 19 & 3.6 {\scriptsize[3.5,4.0]} & 57.7 {\scriptsize[55.4,59.2]} & 0.12 {\scriptsize[0.10,0.13]} & 27.2 & 100/100 \\
dws\_relu\_none\_16\_32\_64 & 65.4 & 20 & 16.1 {\scriptsize[16.0,16.5]} & 203.1 {\scriptsize[193.2,205.1]} & 0.22 {\scriptsize[0.22,0.22]} & 28.9 & 100/100 \\
dws\_relu\_none\_32\_64\_128 & 71.3 & 21 & 30.7 {\scriptsize[30.2,31.0]} & 360.8 {\scriptsize[354.2,373.1]} & 0.43 {\scriptsize[0.42,0.48]} & 29.0 & 100/100 \\
dws\_relu\_none\_32\_64\_128\_256 & 75.8 & 22 & 46.2 {\scriptsize[45.4,49.1]} & 749.9 & 1.56 & 29.0 & 100/100 \\
cnn\_relu\_none\_16\_32\_64 & 77.9 & 22 & 41.0 {\scriptsize[40.4,42.1]} & 712.8 & 1.16 & 29.0 & 100/100 \\
cnn\_relu\_none\_16\_32\_64\_128 & 82.5 & 22 & 42.6 {\scriptsize[41.7,44.1]} & 970.1 & 2.06 & 29.0 & 100/100 \\
\bottomrule
\end{tabular}
\end{table}


\subsection{Budget-Constrained Architecture Search}

\textbf{Protocol.}  The pool is the 243 calibrated trained networks; for a budget $B$ it contains those with calibrated $\logrows\le B$.  A method selects $K=10$ networks to train, and the outcome is the best mean accuracy among them.
Each of 300 random $50/50$ splits of the pool tunes $\lambda$ for the joint proxy on one half and evaluates all methods on the other.  Table~\ref{tab:search} reports regret against the best network in the evaluation half.
Failed networks are not in the pool, so no method can be penalized for selecting one; this is a limitation.

\textbf{Result.}  Under calibrated budgets the tuned joint proxy is indistinguishable from NEAR-only: $\lambda=0$ is chosen in 95\% of tuning splits, so the two select the same networks, and Joint ties NEAR-only in almost
every split at $K=10$ (win/tie/loss against NEAR-only: 0/300/0, 0/300/0, 0/300/0, 19/280/1 at $B=20,21,22,23$).  We therefore withdraw the earlier claim that the joint proxy improves on NEAR-only.
NEAR-only does beat random selection by a wide margin at every budget (regret 0.02 vs.\ 1.58 at $B=20$, 0.04 vs.\ 4.18 at $21$, 0.33 vs.\ 7.47 at $22$ and
3.61 vs.\ 8.63 at $23$ accuracy points), and ranking by cost alone is poor (1.77 to 19.35).  The pools are mixed: ReLU networks are 28--33\% of each, and among the $K=10$ highest-NEAR networks at $B=20,\dots,23$ there are 7, 7, 7 and 3 polynomial
networks.  Because $K=10$ is a large share of the smaller pools, we repeat the comparison with $K=5$: regret is then 0.88, 0.39, 4.17 and 7.07 for NEAR-only against 3.43, 6.62, 11.78 and 12.91 for random selection,
and the tuned joint proxy is again no better than NEAR-only (win/tie/loss 4/290/6, 0/294/6, 0/300/0, 2/266/32; slightly worse at $B=23$).

\textbf{Where the cost model helps.}  The cost model is useful when ezkl is not run.  Table~\ref{tab:estbudget} filters by the model's \emph{estimated} $\logrows$:
0--21\% of the networks it admits exceed the calibrated budget, and it rejects 5--25\% of the networks that are feasible.  Choosing high-NEAR networks with no budget
violates the calibrated budget in every case, whereas choosing inside the estimated budget does so for a fraction of 0.0, 0.4, 0.1 and 0.0 of the $K=10$ selections at $B=20,\dots,23$.  The estimate thus removes most
but not all budget violations; it does not replace calibration.

\begin{table}[h]
\centering
\caption{Budget-constrained search under \emph{ezkl-calibrated} budgets ($\logrows\le B$ from calibration of trained
  networks).  Each of 300 random $50/50$ splits of the pool tunes $\lambda\in\{0,0.25,0.5,1,2,4\}$ for the joint
  proxy on one half and evaluates all methods on the other, choosing $K{=}10$ architectures to train; the outcome is the
  best accuracy among them.  Entries are mean regret in accuracy points against the best architecture in the
  evaluation half, with the $2.5$--$97.5$ percentile range over splits; Random averages $50$ draws per split.  Last
  column: splits where Joint beats / ties / loses to NEAR-only (regret).}
\label{tab:search}
\small
\begin{tabular}{rrrrrrrr}
\toprule
\textbf{$\logrows$} & \textbf{pool} & \textbf{oracle} & \textbf{Joint} & \textbf{NEAR-only} & \textbf{ZK-cost-only} & \textbf{Random} & \textbf{J v N} \\
\midrule
$\le20$ & 36 & 65.4 & 0.02 {\scriptsize[0.00,0.00]} & 0.02 {\scriptsize[0.00,0.00]} & 1.77 {\scriptsize[0.00,5.17]} & 1.58 {\scriptsize[0.44,3.15]} & 0/300/0 \\
$\le21$ & 72 & 71.3 & 0.04 {\scriptsize[0.00,0.00]} & 0.04 {\scriptsize[0.00,0.00]} & 7.10 {\scriptsize[0.45,13.86]} & 4.18 {\scriptsize[1.59,6.96]} & 0/300/0 \\
$\le22$ & 111 & 82.5 & 0.33 {\scriptsize[0.00,3.50]} & 0.33 {\scriptsize[0.00,3.50]} & 19.07 {\scriptsize[10.59,25.09]} & 7.47 {\scriptsize[4.53,11.10]} & 0/300/0 \\
$\le23$ & 175 & 82.5 & 3.39 {\scriptsize[0.00,7.34]} & 3.61 {\scriptsize[0.00,7.34]} & 19.35 {\scriptsize[11.50,26.02]} & 8.63 {\scriptsize[4.97,11.92]} & 19/280/1 \\
\bottomrule
\end{tabular}
\end{table}

\begin{table}[h]
\centering
\caption{Filtering by the cost model's \emph{estimated} $\logrows$ (no ezkl run) against the \emph{calibrated} budget, on the 234 calibrated trained networks.
  ``feasible'': calibrated $\logrows\le B$.  ``admitted'': estimated $\logrows\le B$.  ``Exceeds'': share of admitted networks whose calibrated $\logrows>B$.
  ``Rejected'': share of feasible networks the estimate rejects.  ``Violation'': share of the $K{=}10$ highest-NEAR networks, chosen inside the estimated budget
  (or with no budget), whose calibrated $\logrows>B$.  $95\%$ bootstrap intervals over networks.}
\label{tab:estbudget}
\small
\begin{tabular}{rrrrrrr}
\toprule
$B$ & \textbf{feasible} & \textbf{admitted} & \textbf{exceeds} & \textbf{rejected} & \textbf{violation (est.\ budget)} & \textbf{violation (no budget)} \\
\midrule
$\le20$ & 36 & 33 & 0.18 {\scriptsize[0.05,0.31]} & 0.25 {\scriptsize[0.11,0.40]} & 0.0 {\scriptsize[0.0,0.2]} & 1.0 \\
$\le21$ & 72 & 72 & 0.21 {\scriptsize[0.12,0.30]} & 0.21 {\scriptsize[0.12,0.29]} & 0.4 {\scriptsize[0.1,0.8]} & 1.0 \\
$\le22$ & 111 & 126 & 0.17 {\scriptsize[0.11,0.25]} & 0.05 {\scriptsize[0.02,0.10]} & 0.1 {\scriptsize[0.0,0.3]} & 1.0 \\
$\le23$ & 175 & 165 & 0.00 {\scriptsize[0.00,0.00]} & 0.06 {\scriptsize[0.02,0.09]} & 0.0 {\scriptsize[0.0,0.0]} & 1.0 \\
\bottomrule
\end{tabular}
\end{table}


\subsection{Cost-Model Fit and Evaluation}

Table~\ref{tab:costfit} fits the coefficients of Eq.~\eqref{eq:cost} by non-negative least squares on 187 trained networks and evaluates on the 56 networks of the earlier validation set (excluded from the fit).
On the evaluation split the hand-set proxy matches the calibrated $\logrows$ exactly for 39 and within one for 52 of 56 (Spearman 0.82, RMSE 0.18 in $\log_2$ cells);
the fitted model for 43 and 53 (0.86, 0.11).  Fitted coefficients are close to one assigned cell per multiply-accumulate (standard convolution 0.93, depthwise
1.11, Gemm 0.99) and \emph{zero} per matrix-multiply or convolution node, so the hand-set per-node overhead of $4096$ cells is not supported by the data.  Per-element costs are
15.04 (ReLU), 12.45 (power), 5.55 (normalization), 2.21 (multiplication) and 3.47 (addition), so ReLU and the polynomial power operator cost the same order of magnitude per element, consistent
with Section~\ref{sec:findings}.  Agreement is weaker for $x^3{+}x$ networks (exact on 3/7 of the evaluation networks), where calibration, not cell count, sets the circuit size.  $\varepsilon$ in Eq.~\eqref{eq:cost} is reported
only through these residuals, and we do not claim an exact formula.

\begin{table}[h]
\centering
\caption{Cost model of Eq.~\eqref{eq:cost} against ezkl-calibrated $\logrows$ on \emph{trained} networks with real inputs.
  Coefficients are fitted by non-negative least squares (relative-error weighting) on 180 networks and evaluated on the
  54 networks of the earlier validation set, which are excluded from the fit.  ``exact'' and ``within 1'' count
  networks whose predicted $\lceil\log_2 C\rceil$ equals or is within one of the calibrated $\logrows$; $\rho$ is Spearman between predicted and
  calibrated $\logrows$; RMSE is in $\log_2$ of assigned cells.}
\label{tab:costfit}
\small
\begin{tabular}{lrrrrr}
\toprule
 & $n$ & \textbf{exact} & \textbf{within 1} & $\rho$ & \textbf{RMSE} \\
\midrule
Hand-set proxy, evaluation split & 54 & 37 & 50 & 0.80 & 0.19 \\
Fitted (NNLS), evaluation split & 54 & 39 & 51 & 0.84 & 0.11 \\
Fitted (NNLS), fit split & 180 & 125 & 171 & 0.89 & 0.11 \\
\quad fitted, eval, ReLU & 20 & 16 & 20 & 0.94 & 0.10 \\
\quad fitted, eval, $x^2$ & 11 & 9 & 11 & 0.96 & 0.07 \\
\quad fitted, eval, $x^2{+}x$ & 16 & 11 & 16 & 0.94 & 0.15 \\
\quad fitted, eval, $x^3{+}x$ & 7 & 3 & 4 & 0.12 & 0.09 \\
5-fold CV, all 234 networks: hand-set & 234 & 68\% & 93\% & 0.87 & -- \\
5-fold CV, all 234 networks: fitted & 234 & 70\% & 95\% & 0.89 & -- \\
\quad CV fitted, resnet & 16 & 100\% & -- & -- & -- \\
\quad CV fitted, dws-resnet & 66 & 66\% & -- & -- & -- \\
\quad CV fitted, mlp & 52 & 83\% & -- & -- & -- \\
\quad CV fitted, cnn & 53 & 76\% & -- & -- & -- \\
\quad CV fitted, dws-cnn & 47 & 44\% & -- & -- & -- \\
\bottomrule
\end{tabular}

\smallskip
{\scriptsize Fitted coefficients (cells per MAC / per node / per element): macs\_conv\_standard=0.941, macs\_conv\_depthwise=1.11, macs\_gemm=0.992, n\_matmul\_conv\_nodes=0, elem\_relu=14.8, elem\_pow=12.3, elem\_mul=2.28, elem\_add\_sub=3.4, elem\_norm=5.57, elem\_pool\_reduce=0.}
\end{table}



